
\begin{center}
  {\Large\itshape Alban (au temps des Croisades : XII\textsuperscript{ème} siècle)}
\end{center}

\vspace{1em}

Le 5 mai 1688, le visiteur de l'ordre des hospitaliers de St Jean de Jérusalem fait démettre Jean Rollandet, curé de Montfranc de ses droits, sur les biens des hospitaliers du Carmé. En 1767, le commandeur de Rayssac perçoit encore des redevances à Margou. Le 4 septembre 1778, Guillaume Louis Auguste Passemar, de Capriol Saint Maurice, sieur de Cantelosa, Baron d'Alban, vicomte de St André, affirme à Jean Barry le domaine de la « Mare » qui appartient à son épouse, Rose de Clergue.

Les hospitaliers s'implantent donc dans la région d'assez bonne heure et au cours des siècles, ils auront beaucoup de biens fonciers dans le coin : notez Alban, Montfranc, Rayssac, Carmenel, la Maie, la Font de Carmé, le mas de Carmé, la Boyssounade. Et si « Carmenel » signifie « le petit Carmé », on voit combien ont prospéré des hospitaliers établis d'abord à Montfranc (St Léonard du Carmé) et essaimant dans la région.

Notons que le mot « maio » qui semble bien signifier maison ou communauté de frères hospitaliers, s'explique très bien au voisinage de la font (de Carmé) et du mas de Carmé, ce dernier disparu.

Par la suite, une dame ou demoiselle a fait graver son écu avec une maille (ou fils entrecroisés) sur une pierre avec en accolade -- qui sert de margelle à la font de la maye -- Est-ce cette Rose de Clergue qui devint l'épouse du Passemar, sieur de Cantelauze, baron d'Alban -- ci-dessus nommé ?


